\documentclass[a4paper,11pt]{ltjsarticle}
\usepackage{luatexja}
\usepackage{luatexja-fontspec}
\usepackage{amsmath,amssymb,amsthm}
\usepackage{mathtools}
\usepackage{booktabs}
\usepackage{longtable}
\usepackage{array}
\usepackage{hyperref}
\usepackage{graphicx}
\usepackage{float}
\usepackage{geometry}
\geometry{margin=25mm}

\hypersetup{
  colorlinks=true,
  linkcolor=blue,
  urlcolor=blue,
  citecolor=blue
}

\theoremstyle{definition}
\newtheorem{theorem}{定理}[section]
\newtheorem{proposition}[theorem]{命題}
\newtheorem{lemma}[theorem]{補題}
\newtheorem{corollary}[theorem]{系}
\newtheorem{definition}[theorem]{定義}
\newtheorem{remark}[theorem]{注意}
\newtheorem{observation}[theorem]{観察}

\setlength{\parindent}{1em}
\setlength{\parskip}{0.5em}

\providecommand{\tightlist}{\setlength{\itemsep}{0pt}\setlength{\parskip}{0pt}}
\begin{document}

\section{中心投影フレームワークにおける Schwarzschild--de Sitter
厳密解}\label{ux4e2dux5fc3ux6295ux5f71ux30d5ux30ecux30fcux30e0ux30efux30fcux30afux306bux304aux3051ux308b-schwarzschildde-sitter-ux53b3ux5bc6ux89e3}

\textbf{--- 論文 {[}1{]} の真空 Einstein 方程式の球対称静的厳密解 ---}

\begin{center}\rule{0.5\linewidth}{0.5pt}\end{center}

\textbf{著者}: 木原 範昭 (WF System Co., Ltd.)

\textbf{日付}: 2026年4月

\textbf{DOI：}
\href{https://doi.org/10.5281/zenodo.19538098}{10.5281/zenodo.19538098}

\begin{center}\rule{0.5\linewidth}{0.5pt}\end{center}

\subsection{要旨}\label{ux8981ux65e8}

本稿では、先行論文 {[}1{]} で導出された真空 Einstein 方程式
\(G_{\mu\nu} + \Lambda\, g_{\mu\nu} = 0\)（\(\Lambda = 3/R^2\)）の球対称静的厳密解として、Schwarzschild--de
Sitter 計量を構成する。この厳密解は論文 {[}1{]}
の方程式から追加の仮定なしに一意に決定される（Birkhoff
の定理による）。構成された厳密解において \(R \to 0\) の極限をとり、論文
{[}4{]}
で指摘された構造的緊張がこの厳密解の言葉でどのように表現されるかを記述する。

\textbf{明記事項}: 本稿は論文 {[}1{]}
で導出された方程式の数学的厳密解を構成するのみであり、この厳密解が現実の天体物理学的対象を記述するとは主張しない。Schwarzschild--de
Sitter
計量は一般相対論でよく知られた解であり、本稿はその解を中心投影フレームワークの文脈に位置付ける。物理的解釈は本稿の対象外であり、読者に委ねる。

\begin{center}\rule{0.5\linewidth}{0.5pt}\end{center}

\subsection{1. 導入}\label{ux5c0eux5165}

先行論文 {[}1{]} 命題 6.1 において、中心投影 \(\Phi : \Pi_R \to S^4(R)\)
による誘導計量 \(g_{\mu\nu}\) に対して、真空 Einstein 方程式

\[G_{\mu\nu} + \Lambda\, g_{\mu\nu} = 0, \qquad \Lambda = \frac{3}{R^2} \tag{1.1}\]

が導出された。この方程式は幾何学的恒等式として証明されたものであり（論文
{[}1{]} 注意 6.1）、追加の物理的仮定を含まない。

本稿の目的は、方程式 (1.1) の球対称静的厳密解を構成し、曲率半径 \(R\)
のパラメータ依存性、特に \(R \to 0\)
の極限における振る舞いを検証することにある。

\begin{center}\rule{0.5\linewidth}{0.5pt}\end{center}

\subsection{2.
前提条件の確認}\label{ux524dux63d0ux6761ux4ef6ux306eux78baux8a8d}

方程式 (1.1)
の球対称静的厳密解を求めるにあたり、中心投影フレームワークが持つ構造的性質を確認する。

\subsubsection{2.1 球対称性}\label{ux7403ux5bfeux79f0ux6027}

中心投影 \(\Phi\) は、背景空間 \(\mathbb{R}^{n+1}\) における半径 \(R\)
の超球面 \(S^n(R)\) への射影として定義される（論文 {[}1{]}
§2）。この構成は中心軸の周りの回転に対して不変であり、球対称性を持つ。

\subsubsection{2.2 真空条件}\label{ux771fux7a7aux6761ux4ef6}

方程式 (1.1) の右辺はゼロである。すなわち、エネルギー・運動量テンソル
\(T_{\mu\nu}\) の寄与を含まない真空方程式である。

\subsubsection{\texorpdfstring{2.3 宇宙項 \(\Lambda\)
の由来}{2.3 宇宙項 \textbackslash Lambda の由来}}\label{ux5b87ux5b99ux9805-lambda-ux306eux7531ux6765}

\(\Lambda = 3/R^2\)
は、中心投影の幾何学的構造から導出された量である（論文 {[}1{]} 式
6.2）。一般相対論において外部から導入される宇宙定数とは異なり、本フレームワークでは
\(\Lambda\) は曲率半径 \(R\) の関数として\textbf{内在的に決定される}。

\subsubsection{2.4 定常性}\label{ux5b9aux5e38ux6027}

論文 {[}1{]} の構成において、主観空間 \(\Pi_R\)
は位置情報および時間情報を含まない。誘導計量 \(g_{\mu\nu}\) は座標
\(x^\mu\)
にのみ依存し、時間変数を含まない。この性質は、厳密解が静的であることと整合する。

\begin{center}\rule{0.5\linewidth}{0.5pt}\end{center}

\subsection{3. Schwarzschild--de Sitter
厳密解の構成}\label{schwarzschildde-sitter-ux53b3ux5bc6ux89e3ux306eux69cbux6210}

\subsubsection{3.1 Birkhoff
の定理の適用}\label{birkhoff-ux306eux5b9aux7406ux306eux9069ux7528}

一般相対論における Birkhoff の定理の \(\Lambda\) 付き拡張（Kottler,
1918）により、球対称かつ真空の Einstein 方程式

\[G_{\mu\nu} + \Lambda\, g_{\mu\nu} = 0 \tag{3.1}\]

の解は、静的であり、かつ以下の Schwarzschild--de Sitter 計量（Kottler
計量）に一意に帰着する:

\[ds^2 = -f(r)\,dt^2 + f(r)^{-1}\,dr^2 + r^2\,d\Omega^2 \tag{3.2}\]

ここで

\[f(r) = 1 - \frac{2M}{r} - \frac{\Lambda}{3}\,r^2 \tag{3.3}\]

であり、\(M\) は積分定数（質量パラメータ）、\(d\Omega^2\)
は単位球面上の計量である。

\subsubsection{3.2
中心投影フレームワークへの代入}\label{ux4e2dux5fc3ux6295ux5f71ux30d5ux30ecux30fcux30e0ux30efux30fcux30afux3078ux306eux4ee3ux5165}

方程式 (1.1) の \(\Lambda = 3/R^2\) を式 (3.3) に代入すると

\[f(r) = 1 - \frac{2M}{r} - \frac{r^2}{R^2} \tag{3.4}\]

を得る。これが中心投影フレームワークにおける Schwarzschild--de Sitter
厳密解である。

\subsubsection{3.3 特殊な場合}\label{ux7279ux6b8aux306aux5834ux5408}

\textbf{\(M = 0\) の場合}: 式 (3.4) は

\[f(r) = 1 - \frac{r^2}{R^2} \tag{3.5}\]

となり、純粋な de Sitter
計量に帰着する。これは質量源のない場合であり、中心投影の誘導計量そのものの球対称表示に対応する。

\textbf{\(R \to \infty\)（\(\Lambda \to 0\)）の場合}: 式 (3.4) は

\[f(r) \to 1 - \frac{2M}{r} \tag{3.6}\]

となり、標準的な Schwarzschild 計量に帰着する。これは論文 {[}4{]} §2 の
\(R \to \infty\) 極限と整合する。

\begin{center}\rule{0.5\linewidth}{0.5pt}\end{center}

\subsection{\texorpdfstring{4. \(R \to 0\)
極限の検証}{4. R \textbackslash to 0 極限の検証}}\label{r-to-0-ux6975ux9650ux306eux691cux8a3c}

\subsubsection{\texorpdfstring{4.1 \(\Lambda\)
の発散}{4.1 \textbackslash Lambda の発散}}\label{lambda-ux306eux767aux6563}

\(R \to 0\) の極限において \(\Lambda = 3/R^2 \to \infty\) となる。式
(3.4) において \(r^2/R^2\) の項が支配的になり

\[f(r) \approx -\frac{r^2}{R^2} \tag{4.1}\]

が \(r > 0\) の任意の値で成立する。\(f(r) < 0\)
であるから、\(r = \text{const}\)
の面は時間的（すなわち地平面の内側の構造に対応する）となる。

\subsubsection{4.2
地平面の構造}\label{ux5730ux5e73ux9762ux306eux69cbux9020}

\(f(r) = 0\) の解は地平面の位置を与える。式 (3.4) から

\[1 - \frac{2M}{r} - \frac{r^2}{R^2} = 0 \tag{4.2}\]

\(R\) が十分大きい場合、二つの正の解が存在する:

\begin{itemize}
\tightlist
\item
  \textbf{ブラックホール地平面} \(r_{\mathrm{BH}}\): \(r \approx 2M\)
  付近
\item
  \textbf{宇宙論的地平面} \(r_{\mathrm{C}}\): \(r \approx R\) 付近
\end{itemize}

\(R \to 0\) の極限では \(r_{\mathrm{C}} \to 0\)
となり、宇宙論的地平面がブラックホール地平面に接近して合流する。すなわち、二つの地平面の間の「観測者が存在しうる領域」が消失する。

\subsubsection{4.3 論文 {[}4{]}
の構造的緊張との対応}\label{ux8ad6ux6587-4-ux306eux69cbux9020ux7684ux7dcaux5f35ux3068ux306eux5bfeux5fdc}

論文 {[}4{]} §3.2 で指摘された構造的緊張 --- 測地線偏差の上限
\(|\xi|^2_{\max}(R)\) がゼロに漸近する一方で局所曲率が発散する ---
は、本節の厳密解の言葉で次のように表現される:

\(R \to 0\) において宇宙論的地平面 \(r_{\mathrm{C}} \to 0\)
となり、内部観測者が存在しうる領域が消失する。これは論文 {[}4{]} §3.2
の「主観空間内の観測者が測定可能な量がゼロに漸近する」という記述の厳密解レベルでの対応物である。

\subsubsection{4.4
本節の主張の範囲}\label{ux672cux7bc0ux306eux4e3bux5f35ux306eux7bc4ux56f2}

本節で述べたのは以下に限定される。

\begin{enumerate}
\def\labelenumi{\arabic{enumi}.}
\tightlist
\item
  論文 {[}1{]} の方程式 (1.1) の厳密解 (3.4) において、\(R \to 0\)
  の極限が幾何学的にどのような構造を示すか。
\item
  この構造が論文 {[}4{]} §3.2 の構造的緊張とどのように対応するか。
\end{enumerate}

本稿は、この厳密解が特定の天体物理学的対象（ブラックホールなど）を記述するとは主張しない。\(M\)
は数学的な積分定数であり、物理的質量との同一視は本稿の射程を超える。

\begin{center}\rule{0.5\linewidth}{0.5pt}\end{center}

\subsection{5.
入れ子構造との接続}\label{ux5165ux308cux5b50ux69cbux9020ux3068ux306eux63a5ux7d9a}

\subsubsection{5.1 命題 4.1
の適用}\label{ux547dux984c-4.1-ux306eux9069ux7528}

§4 で記述した \(R \to 0\)
極限において内部観測者の領域が消失する状況に対し、論文 {[}4{]} 命題
4.1（入れ子構造の存在）が適用可能である。\(R > 0\)
の任意の値に対して、主観空間 \(\Pi_R\) の内部の任意の点から、曲率半径
\(R' > 0\) を持つ新たな主観空間 \(\Pi_{R'}\) が構成可能であり、\(R'\) は
\(R\) とは独立に選べる。

\subsubsection{5.2
子主観空間における厳密解}\label{ux5b50ux4e3bux89b3ux7a7aux9593ux306bux304aux3051ux308bux53b3ux5bc6ux89e3}

子主観空間 \(\Pi_{R'}\) においても、論文 {[}1{]}
の構成がそのまま適用される。したがって、\(\Pi_{R'}\) における真空
Einstein 方程式は

\[G'_{\mu\nu} + \Lambda'\, g'_{\mu\nu} = 0, \qquad \Lambda' = \frac{3}{R'^2} \tag{5.1}\]

であり、その球対称静的厳密解は

\[f'(r') = 1 - \frac{2M'}{r'} - \frac{r'^2}{R'^2} \tag{5.2}\]

である。\(R'\) は \(R\) とは独立であるから、親主観空間で \(R \to 0\)
であっても、子主観空間では \(R'\) を有限に保つことができる。

\subsubsection{5.3
地平面構造の回復}\label{ux5730ux5e73ux9762ux69cbux9020ux306eux56deux5fa9}

親主観空間において §4.2
で消失した「観測者が存在しうる領域」は、子主観空間 \(\Pi_{R'}\)
の内部では \(R' > 0\) が有限であるかぎり回復する。子主観空間の厳密解
(5.2) は、\(R'\)
の値に応じたブラックホール地平面と宇宙論的地平面を持ち、両地平面の間に観測者が存在しうる
well-defined な領域が存在する。

\subsubsection{5.4
本節の主張の範囲}\label{ux672cux7bc0ux306eux4e3bux5f35ux306eux7bc4ux56f2-1}

本節で述べたのは以下に限定される。

\begin{enumerate}
\def\labelenumi{\arabic{enumi}.}
\tightlist
\item
  論文 {[}4{]} 命題 4.1 により、\(R \to 0\)
  の極限領域から子主観空間が構成可能であること。
\item
  子主観空間においても論文 {[}1{]} の方程式とその厳密解が成立すること。
\item
  親主観空間で消失した地平面間領域が、子主観空間の内部では回復すること。
\end{enumerate}

これらは論文 {[}1{]} 命題 6.1 と論文 {[}4{]} 命題 4.1
の直接の帰結であり、新たな仮定を含まない。物理的解釈は本稿の対象外であり、読者に委ねる。

\begin{center}\rule{0.5\linewidth}{0.5pt}\end{center}

\subsection{6. 考察}\label{ux8003ux5bdf}

本稿では、論文 {[}1{]} 命題 6.1 で導出された真空 Einstein
方程式の球対称静的厳密解として Schwarzschild--de Sitter
計量を構成し、\(R \to 0\) 極限における振る舞いを論文 {[}4{]}
の入れ子構造と接続した。

本稿の記述は以下に限定される。

\begin{enumerate}
\def\labelenumi{\arabic{enumi}.}
\tightlist
\item
  論文 {[}1{]} の方程式 (1.1) と Birkhoff の定理から、Schwarzschild--de
  Sitter 解が追加の仮定なしに一意に決定されること。
\item
  \(\Lambda = 3/R^2\) の代入により、厳密解が \(R\)
  のパラメータを持つ具体的な形で記述されること。
\item
  \(R \to 0\) の極限において内部観測者の領域が消失し、これが論文 {[}4{]}
  §3.2 の構造的緊張の厳密解レベルでの対応物であること。
\item
  論文 {[}4{]} 命題 4.1
  の入れ子構造により、子主観空間において地平面間領域が回復すること。
\end{enumerate}

本稿は、この厳密解が特定の物理的状況を記述するとは主張しない。\(M\)
は積分定数であり、\(R\)
は中心投影の曲率半径であり、いずれも物理的量との同一視は本稿の射程を超える。Schwarzschild--de
Sitter
計量に関連する物理学的概念（ブラックホール、宇宙論的地平面、Hawking
輻射、情報パラドックスなど）との対応は、幾何学の枠組みの外側の問題であり、本稿では扱わない。

\subsubsection{シリーズの論理的読み順について}\label{ux30b7ux30eaux30fcux30baux306eux8ad6ux7406ux7684ux8aadux307fux9806ux306bux3064ux3044ux3066}

本シリーズの論文は公開順に番号付けされているが、論理的な読み順としては以下を推奨する:

\begin{itemize}
\tightlist
\item
  論文 {[}1{]}--{[}3{]}: 基礎フレームワーク
\item
  論文 {[}4{]}: \(R\) 極限と入れ子構造
\item
  論文 {[}5{]}: 次元追加
\item
  \textbf{本稿（論文 {[}9{]}）: 厳密解と \(R \to 0\) の地平面構造}
\item
  論文 {[}7{]}: 測地線の橋渡し
\item
  論文 {[}6{]}: 各次元の主観空間の記述
\item
  論文 {[}8{]}: 0次元からの次元展開
\end{itemize}

\begin{center}\rule{0.5\linewidth}{0.5pt}\end{center}

\subsection{7. 結論}\label{ux7d50ux8ad6}

中心投影フレームワークの真空 Einstein 方程式
\(G_{\mu\nu} + \Lambda g_{\mu\nu} = 0\)（\(\Lambda = 3/R^2\)）の球対称静的厳密解は、Birkhoff
の定理（の \(\Lambda\) 付き拡張）により、Schwarzschild--de Sitter 計量
\(f(r) = 1 - 2M/r - r^2/R^2\)
に一意に決定される。この構成に追加の仮定は不要である。\(R \to 0\)
の極限において内部観測者の領域は消失するが、論文 {[}4{]} 命題 4.1
の入れ子構造により、子主観空間において同じ厳密解の構造が \(R'\)
を独立なパラメータとして回復する。物理的解釈は本稿の対象外であり、読者に委ねる。

\begin{center}\rule{0.5\linewidth}{0.5pt}\end{center}

\subsection{参考文献}\label{ux53c2ux8003ux6587ux732e}

{[}1{]} 木原 範昭,「中心投影による4次元空間の幾何学的定式化」, Zenodo,
DOI: 10.5281/zenodo.19427780 (2026).

{[}2{]} 木原 範昭,「中心投影の幾何学的対称性：多軸モデルの数学的基盤」,
Zenodo, DOI: 10.5281/zenodo.19434932 (2026).

{[}3{]} 木原
範昭,「複数の主観空間における観測の相対性：中心投影の対称性の幾何学的帰結」,
Zenodo, DOI: 10.5281/zenodo.19435162 (2026).

{[}4{]} 木原 範昭,「中心投影における曲率半径の極限についての考察 ---
\(R \to \infty\) と \(R \to 0\) の場合」, Zenodo, DOI:
10.5281/zenodo.19526549 (2026).

{[}5{]} 木原 範昭,「背景空間と主観空間への次元追加についての考察」,
Zenodo, DOI: 10.5281/zenodo.19526913 (2026).

{[}6{]} 木原 範昭,「ゼロ次元から四次元主観空間の考察」, Zenodo, DOI:
10.5281/zenodo.19533292 (2026).

{[}7{]} 木原 範昭,「主観空間における測地線の連続性についての考察 ---
親主観空間から子主観空間への測地線の橋渡しについて」, Zenodo, DOI:
10.5281/zenodo.19533299 (2026).

{[}8{]} 木原 範昭,「体積1の四次元超直方体に外接する超球体の直径」,
Zenodo, DOI: 10.5281/zenodo.19533313 (2026).

\begin{center}\rule{0.5\linewidth}{0.5pt}\end{center}

\end{document}
